\chapter*{前言}
\markboth{前言}{前言}
\addcontentsline{toc}{chapter}{前言}

\textbf{\trFukayaCategory}（Fukaya category）这门学科素有难以入门之名。这既是由于所需的背景知识繁多（取自\trHomology{}代数、\trSymplectic{}几何与几何分析），同样也是由于基本定义本身就相当繁复。本书旨在为想学习\trFukayaCategory、并可能在自己工作中加以运用的研究生与研究人员提供一份资源。我力求聚焦于相当基础的一批论题，并尽可能精确地描述它们，同时补上早期参考文献中的若干缺漏。这造就了一种相当朴素冷峻的文体（因此，要想透彻研读本书，恐怕还应辅以阅读若干讨论其应用的论文）。第二个目的是记述一些此前未曾发表的结果，它们将\trFukayaCategory{}与\textbf{\trLefschetzFibration}（Lefschetz fibration）的理论联系起来。这类内容在最后几节中占据主导，全书在那里逐渐转为一部研究专著。

我从许多人的工作中大量取材，其中首推\trFukaya、\trKontsevich{}与\trDonaldson。\trFukaya{}的奠基性贡献，当然是将\trAinfty{}结构引入\trSymplectic{}几何。在代数方面，他开创了\trYonedaEmbedding{}的\trAinfty{}版本的使用，我们系统地采纳了这一做法。除此之外，若干几何想法——例如\trPinStructure{}的作用，以及用\trParametrization{}的\trModuliSpace{}构造\trAinfty-\trHomomorphism——都取自\trFukaya、\trOh、\trOhta{}与\trOno{}的工作。\trKontsevich{}引入了\textbf{\trACategory}（$A_\infty$-category）的\trDerivedCategory，这些\trCategory{}理论的许多部分都出自他手，特别是\textbf{\trExactTriangle}（exact triangle）的内在刻画。他还猜想\textbf{\trDehnTwist}（Dehn twist）与\trTwistFunctor{}之间的联系，这是本书的主要结果之一。最后，在与\trBarannikov{}的合作中，他提出了为\trLefschetzFibration{}构造\trFukayaCategory{}的一种方案；我们采用的定义表面上不同，但想必是\trEquivalence{}的。\trDonaldson{}的影响同样无处不在。除了他在\trLefschetzPencil{}方面的开创性工作之外，他还引入了\trMatchingCycle，并提议以其为\trFloerCohomology{}某个组合公式的出发点；这一设想在本书中确有部分实现。其他数学家也做出过重要贡献。例如，我们对\textbf{\trPicardLefschetzTheory}（Picard--Lefschetz theory）的表述就有一部分反映了\trAuroux{}的观点。\trSmith{}的一个结果，即四维\trLefschetzPencil{}中的\trVanishingCycle（vanishing cycle）必定填满整条\trFiber，对于暗示这类\trCycle{}或许能“\trSplitGeneration”（split-generate）\trFukayaCategory{}起了关键作用。此外，\trFukaya--\trSmith{}关于\trCotangent{}丛的工作，为关于\trLefschetzFibration{}的若干更大胆的想法提供了良好的试验场。我们处理\trTransversality{}问题的方法，来自与\trLazzarini{}的多次交谈。最后，\trAbouzaid{}的建议大大改进了关于\trSymplectic\trEmbedding{}的讨论。

在机构方面，我要感谢ETH Zürich邀请我讲授本书所依据的那门课程；感谢NSF的经费支持（资助编号DMS-9729992与DMS-0405516）；也感谢EMS Publishing House接受书稿，并护送其顺利出版。
